% Generated draft: requires mathematical and source review.
\documentclass[11pt]{article}
\usepackage[margin=25mm]{geometry}
\usepackage{fontspec}
\setmainfont{Latin Modern Roman}
\usepackage{amsmath,amssymb,mathtools,unicode-math}
\setmathfont{Latin Modern Math}
\usepackage{enumitem,needspace,xurl,hyperref}
\setlength{\parindent}{0pt}
\setlength{\parskip}{5pt}
\setlength{\emergencystretch}{3em}
\newcommand{\N}{\mathbb{N}}
\newcommand{\Z}{\mathbb{Z}}
\newcommand{\Q}{\mathbb{Q}}
\newcommand{\R}{\mathbb{R}}
\newcommand{\C}{\mathbb{C}}
\newcommand{\F}{\mathbb{F}}
\newcommand{\A}{\mathbb{A}}
\newcommand{\PP}{\mathbb{P}}
\newcommand{\E}{\mathbb{E}}
\newcommand{\eps}{\varepsilon}
\newcommand{\dd}{\,\mathrm{d}}
\begin{document}
\section*{UnsolvedMath ID: 20001279}
\textbf{Source problem code:} catalogue entry
\subsection{Definitions and mathematical statement}
This extraction merges twelve distinct questions. Use $\operatorname{vol}_n$ for Lebesgue volume, $B_p^n$ for the usual $\ell_p$ unit ball when $p>0$, and, for a convex body containing the origin in its interior,
\[
 K^*=\{y:\langle x,y\rangle\leq1\text{ for every }x\in K\},\qquad
 P(K)=\operatorname{vol}_n(K)\operatorname{vol}_n(K^*).
\]
The standard origin-symmetric interpretation is used where symmetry is intended, but several questions do not state it explicitly. For $n\geq2$, the intersection body $IK$ has radial function $\rho_{IK}(u)=\operatorname{vol}_{n-1}(K\cap u^\perp)$, while the projection body $\Pi K$ has support function $h_{\Pi K}(u)=\operatorname{vol}_{n-1}(K|u^\perp)$. These different functions must not be interchanged. Let $\operatorname{St}_u(K)$ denote Steiner symmetrization along $u$, centering each chord parallel to $u$ about $u^\perp$.
\begin{enumerate}
\item Does $IK\subseteq IL$ together with $\Pi K\subseteq\Pi L$ imply $K\subseteq L$?
\item Under the same assumptions, does $\operatorname{vol}_n(K)\leq\operatorname{vol}_n(L)$? The source reports this true for $n\leq4$; its precise body assumptions are not restated.
\item For every convex body $K$, in particular in the symmetric version, does there exist $u\in S^{n-1}$ with $\operatorname{vol}_n(\operatorname{St}_u(K)^*)\leq C\operatorname{vol}_n(K^*)$? An absolute, dimension-independent $C$ is the natural intended reading but is not explicitly quantified in the extract.
\item Does the corresponding average bound hold? With probability surface measure $\sigma$, this reads $\int_{S^{n-1}}\operatorname{vol}_n(\operatorname{St}_u(K)^*)\,\dd\sigma(u)\leq C\operatorname{vol}_n(K^*)$. The source writes $\dd u$ without normalization, which must be fixed.
\item Is the average volume product of random symmetric polytopes at least $P(B_\infty^n)=4^n/n!$? The source first says at least, then writes a strict $\E P(K)>P(B_\infty^n)$. A natural Gaussian model is $K=\operatorname{conv}\{\pm g_1,\ldots,\pm g_N\}$ for independent standard Gaussian $g_i$ and $N\geq n$, but this fixes generators, not necessarily the actual number of vertices. It also does not make all antipodal vertices independent. Those distinctions and strictness remain unresolved.
\item For an even log-concave probability measure $\mu$, bound $\mu(K)\mu(K^*)$ from above. Even means invariant under $x\mapsto-x$; in a density interpretation log-concavity means a density $e^{-\phi}$ with $\phi$ convex. The source parenthesis literally reports $\mu(K)\mu(K^*)\leq\mu(B_2^n)$ when $\mu$ is unconditional, meaning invariant under coordinate sign changes. The unsquared right side is preserved, and the intended assumptions on $K$ need checking.
\item Characterize equality in the previous question and find lower bounds in the unconditional case; whether unconditionality concerns only $\mu$ or also $K$ is not specified.
\item Consider the dual measure $\mu^*$ in the sense of Artstein--Milman and ask the analogous questions for the expression literally written $\mu(K)\mu^*(K)$. Both arguments are $K$ in the source. The duality and normalization are not defined; a Legendre-dual density convention would need to be specified before numerical bounds or equality can be formulated. It is not silently changed to $\mu(K)\mu^*(K^*)$.
\item For symmetric convex bodies define $N(K,L)=\min\{m:K\subseteq\bigcup_{i=1}^m(x_i+L)\}$. Are there absolute $c,\delta>0$, uniform in dimension and bodies, such that $N(K,L)\leq N(L^*,cK^*)^\delta$? The damaged bracket is repaired without changing the direction of this covering-duality question.
\item For uniform random vectors $X\in K$, $Y\in K^*$, are the first roughly one hundred moments of $\langle X,Y\rangle$ within a factor $1+\epsilon$ of an appropriate Gaussian comparison? Independence, normalization, the Gaussian reference (scalar or inner product), the meaning of roughly one hundred, and an asymptotic regime for $\epsilon$ are unspecified; no precise universal moment assertion is inferred.
\item Prove the symmetric Mahler lower bound $P(B_X)\geq4^n/n!$ for $n$-dimensional normed spaces embedding in $L_p$, $p<1$, especially convex intersection bodies, which the source associates with $p=-1$. For $0<p<1$ one may mean a linear isometric embedding into a quasi-normed $L_p$ space. Negative-$p$ embedding is a different generalized notion, not an ordinary $L_p$ norm, and its definition is absent.
\item For $1<p<2$, the source asks whether $P(B_X)>P(B_p^n)$ when $X$ embeds in $L_p$ and $\dim X=n$. Under the usual isometric meaning, $X=\ell_p^n$ already gives equality, so the literal strict sign needs qualification or replacement; no non-strict repair is silently adopted.
\end{enumerate}

The attached aid addresses only the first two questions, under origin symmetry. It reports that question 1 is affirmative for $n=2$, using $IK=\Pi K=2RK$ with $R$ a quarter-turn, and false for every $n\geq3$: take $L=B_2^n$ and $K=\operatorname{diag}(R_0,s,\ldots,s)B_2^n$, with $R_0>1$ and $s=R_0^{-1/(n-2)}$. More generally, for origin-centered ellipsoids $E,F$, it records
\[
 IE\subseteq IF\ \Longleftrightarrow\ \Pi E\subseteq\Pi F
 \ \Longleftrightarrow\ (\operatorname{vol}_n(E)/\operatorname{vol}_n(F))F\subseteq E.
\]
These equivalent conditions force $\operatorname{vol}_n(E)\leq\operatorname{vol}_n(F)$, settling question 2 for ellipsoid pairs in all $n\geq2$. They do not settle question 2 for all bodies or questions 3--12. All reported results remain supplied claims requiring independent verification.
\subsection{Short English statement}
The record contains twelve questions about sections, projections, symmetrization, random volume products, measure duality, covering, moments, and Mahler inequalities. All are retained. The supplied aid resolves section/projection containment under origin symmetry and proves only the ellipsoid case of volume comparison; other formulations retain their source ambiguities.
\subsection{Sources}
\begin{itemize}
\item AIM workshop \emph{Mahler's conjecture and duality in convex geometry}, August 9--13, 2010, notes by Jaegil Kim, questions 1--12 merged from the first two pages: \url{https://aimath.org/WWN/mahlerduality/mahlerduality.pdf}.
\item The source names Artstein--Milman for measure duality but gives no full reference or dual-measure definition. No separate citations support its low-dimensional or unconditional parenthetical claims in this record.
\item Local registry \texttt{data/unsolvedmath/problems.json}, permanent ID 20001279, exact source code \texttt{catalogue entry}; canonical location \texttt{aim-convex-geometry-notes.json notes[10]}. Attached attempt 2, partial result, problem status at run partially solved, rigor proved, low novelty confidence; machine-generated, not peer reviewed, requiring independent verification before citation.
\end{itemize}
\end{document}
